Classifier-free guidance (CFG) is known to buy class fidelity at the price of diversity, but on images the price is hard to itemise. We train a small conditional DDPM on a 2D mixture with four classes of three modes each (mode weights 0.5/0.3/0.2, interleaved so that classes overlap), where the true class posterior, modes and weights are known, and sweep the guidance scale $w$ from 0 to 8 over five seeds. We compare with unguided sampling and with a reimplemented low-temperature baseline on the same network. Guidance raises class accuracy from 0.926 to 0.999 at $w=1$ on the overlapping task, but it does so by giving up minor modes: at $w=3$ the minor mode keeps 0.63 of its true mass (unguided 0.96), mode coverage falls from 1.000 to 0.850, and the sliced Wasserstein distance to the true class distribution grows monotonically with $w$ (0.120 at $w=0$, 1.473 at $w=8$). Variance does not collapse monotonically: the within-mode spread ratio is lowest around $w=1$--2 (0.748, 0.752) and then rises, while off-support mass grows at large $w$. Low-temperature sampling also raises class accuracy, but at a similar spread ratio ($\tau=0.8$ versus $w=0.5$) CFG reaches 0.990 class accuracy against 0.966, and at $\tau=0.5$ it loses more modes (coverage 0.700) than CFG at $w=3$ (0.850) while removing more variance. Guidance restricted to the low-noise half of the chain kept all modes in our run. The mode-weight distortion (mode\_tv) did not depend measurably on class overlap; variance loss was larger and sliced-Wasserstein growth smaller on the separated task (post hoc, uncorrected). This is a small CPU study (one architecture, five seeds, one mode layout), and the interval and mechanism findings are single-configuration observations.
